
\chapter{Discussion}\label{sec:6}

This chapter offers a final reflection on the findings of the empirical studies. Section 6.1 examines the results from a variationist perspective, while Section 6.2 considers how the same predictors account for the distribution of different focus-marking strategies in the two languages. Section 6.3 situates French and Spanish within a broader typology of focus realization, and Section 6.4 reflects on synchronic variation as well as potential diachronic change in the structures identified in the corpus. Finally, Section 6.5 outlines the limitations of the study and suggests possible directions for future research.

\section{A variationist approach to subject focus}

In this book, the question of subject focus in French and Spanish was investigated by adopting a quantitative variationist perspective. Within the variationist approach, a variable is defined as a general or abstract feature that is subject to variation, while the actual instantiations of the variable in speech are known as the variants \parencite{meyerhoff2018introducing}. By examining the discourse functions of these structures in context, one can discern whether two or more alternating syntactic structures are indeed competing variants and, therefore, belong to the same envelope of variation. 
\\
To investigate the phenomenon of subject focus, I adopted this perspective by starting from the pragmatic contexts that trigger the use of specific syntactic structures. Although these strategies may involve a change in referent and/or syntactic function, they all serve the same information-structural purpose and can therefore be regarded as optional variants. Put differently, the syntactic structures analysed in this book are treated as pragmatically equivalent realizations of the same underlying variable.\\
Furthermore, these structures exhibit noticeable differences at the syntactic level. For instance, when comparing the \textit{in situ} strategies in both languages to an \textit{il y a} cleft, one can observe that the focused element has changed its syntactic function: while it appears as the subject in the \textit{in situ} structure, in the \textit{il y a} cleft, it functions as the direct object of the sentence.\footnote{The same phenomenon occurs with a \textit{c'est} cleft, where the focalized element appears as the attribute of the copula \textit{être.}} Nevertheless, all the strategies found in the Spanish and French datasets fulfil the same semantic and pragmatic requirement previously illustrated by example (\ref{ex504}) in Chapter 3, which I repeat here: 

\ea \label{ex504} 
\ea (Who ate beans?)
\ex \ [Mary]\\textsc{foc} \normalsize ate beans. \hfill \parencite[7]{velleman2016question}    
\z \z

Following \textcite{hamblin1976questions}, the (implicit or explicit) QUD in (\ref{ex504}a) can be interpreted as a set of propositions, each being the denotation of a congruent answer: \( \{ \text{ATE(BEANS)}(x) \mid x \in \text{ENTITY} \} \). The utterance in (\ref{ex504}b), on the other hand,  identifies one of these propositions and adds it to the CG content, stating that of all the possible alternative entities that could have eaten the beans, Mary did. 
The fact that this semantic and pragmatic requirement can appear as a canonical SVO structure, a cleft sentence, a VS order or a reduced cleft is evidence of the fact that there is variability in subject focus realization, and that subject focus itself can be considered a linguistic variable. However, this does not imply that the differences among these structures are neglected. As \textcite{adli2006french} writes, ``if there is optionality at the level of the syntax module, this does not mean that the resulting variants are in free variation for all other modules of grammar” \parencite[198]{adli2006french}. Rather, the \textit{raison d'être} of such optionality must be individuated and accounted for.\\ 
Indeed, a key tenet of the variationist approach is the idea of ``orderly or structured heterogeneity” \parencite[99-100]{weinreich1968empirical}. Variation is not random, but structured in regular ways along a number of internal linguistic and external social dimensions, including, for instance, linguistic context, speaker position/status, and text type (register, style, genre, etc.). 
Just like two representations of the same phoneme can be motivated by, say, differences in the social class of the speakers, one must be able to motivate the alternation among different syntactic forms that compete to express the same function. \\
In this study, I have looked at the pragmatic and linguistic factors that might explain the variation found among the different syntactic strategies. For instance, in Chapter 4, I investigated the alternation between full and elliptical forms when they are used to express the same communicative function. The results have shown that full forms (e.g., VS, \textit{in situ}, clefts, etc.) are found more frequently when answering QUDs that are not immediate or implicit, while elliptical forms (e.g., reduced clefts and fragments) generally occur in the presence of immediately preceding QUDs. These results, however, are interpreted at a probabilistic level. Notably, one variant appearing with a higher frequency in a specific context does not automatically exclude the possibility that the same form could be used in a different context. In other words, it has been shown that a full cleft may well be a felicitous answer to an immediately preceding QUD, or vice versa, a reduced cleft might occur in a context that involves an implicit QUD.\\
Thus, these cases are significantly different from cases where constraints on a variable are entirely regular and predictable, allowing for the consistent determination of which variant will surface. For example, in many phonetic and phonological variables, linguistic cues — such as syllable position — dictate the variant of a particular phoneme that is used. In such instances, the realizations of that phoneme are determined by its syllabic context. In contrast, with the variation observed in subject focus realization, the same speaker may alternate between variants, using one in some contexts and another in others. Even if there is a higher probability that one variant will be favoured over another, no contextual constraint guarantees that a specific form will be employed 100\% of the time. Nevertheless, it remains possible to predict the likelihood of encountering different forms in various contexts and among different speakers, especially in the presence of specific predictors. 

\section{Predictors of focus-marking strategies}
To investigate what influences the alternation between different word orders in Spanish and different cleft formats in French, I examined the possible impacts of the same predictors on the two languages. In particular, I analysed whether the information status of the subject, the verb type contained in the syntactic structure and exhaustivity would be able to predict the choice of one rather than the other form. The results showed again a scenario where certain forms appear with a higher frequency with specific features rather than others. \\
For instance, it has been demonstrated that unaccusative verbs frequently combine with new subjects, while, at least in French, transitive verbs are often paired with given subjects. Moreover, these characteristics tend to correlate with specific syntactic structures, such as verb-subject (VS) order with unaccusative verbs and new subjects. However, these findings indicate tendencies rather than deterministic correlations. The same features may also appear in different word orders or cleft formats, albeit less frequently. This further supports the notion that there is no strict one-to-one correspondence between form and function, and that various syntactic constructions can often be used interchangeably. Therefore, opting for a \textit{c'est} cleft in contexts where speakers would typically use an \textit{il y a} cleft (e.g., with an unaccusative verb and a new subject) does not result in ungrammaticality, and the choice among different forms is primarily influenced by pragmatic constraints rather than solely by syntactic or semantic requirements.\\
Among the three factors investigated, focus type — specifically exhaustivity — has been shown to differ somewhat from the other two. Certain theories have posited that exhaustivity is inherently a semantic notion \parencite{szabolcsi1981semantics, kiss1998identificational}, which carries quantificational effects on the sentence in question. In other words, while pragmatic constraints do not influence the truth-conditional value of a sentence, exhaustivity does. As shown in Chapter 5, exhaustivity affects syntactic variation in both languages, although this effect is more pronounced in the French dataset. This finding might lead one to argue that the constructions investigated cannot be considered optional variants because they do not satisfy the criterion of semantic equivalence due to their differences at the semantic level.\\
However, such a claim is not supported by the data. In the case of Spanish, the effect of exhaustivity on word order is barely perceptible. 
In fact, mixed-effects modelling indicates that the association between focus type (i.e., exhaustivity) and syntactic structure is not even significant. For the French data, on the contrary, the results reveal a clear pattern linking exhaustive focus with \textit{c'est} clefts and not-exhaustive focus with \textit{il y a} clefts. This could be interpreted as evidence against the interchangeability of the two cleft formats. Given that \textit{c'est} clefts express exhaustive focus while \textit{il y a} clefts express its opposite, one might argue that these constructions cannot be used interchangeably in the same context and do not represent optional variants of the same variable. However, even in the case of French, the data does not support this assertion.\\
First, the results indicate probabilistic tendencies rather than deterministic rules. \textit{Il y a} clefts can also convey exhaustivity, and, albeit less frequently, \textit{c'est} clefts can appear in not-exhaustive contexts. Moreover, as argued in Chapter 5, exhaustivity is not syntactically encoded in either construction but is pragmatically inferred from contextual cues. \\
The evidence for this claim, already supported by other researchers (see, for example, \cite{horn1981exhaustiveness, dufter2009clefting, destruel2018interpretation}), is that both types of clefts can be accompanied by particles that indicate either exhaustivity or its absence. This is visible in examples (\ref{ex505}) and (\ref{ex506}), already illustrated in Chapter 2 and Chapter 5, that I report again here:

\ea \label{ex505} 
 \textit{C'est seulement} {[\textit{moi}]\\textsc{foc}} \normalsize \textit{ensuite qui suis venue frapper pour la prévenir}.\\  
`It is only me, afterwards, who came to warn her.' \hfill (sgs French, 38/187)   
\z

\ea \label{ex506}
\textit{Bon, c'est vrai qu'il y avait} {[\textit{ce gars}]\\textsc{foc}} \normalsize \textit{aussi qui venait là, depuis deux, trois semaines là, qui venait assez régulièrement.}\\
`Well, it is true that there was this guy who, since 2 or 3 weeks, was coming over pretty regularly.' \hfill (sgs French, 77/366)  
\z

Thus, differences that may arise between the two constructions are not encoded in the syntactic structure itself, but rather derived by the context in which the sentence is found. As mentioned above, pragmatic differences should not exclude variants from the envelope of variation; instead, they should be considered possible factors that might explain what determines such variation.\\
The occurrence of competing variants for subject focus expression raises the important question of where to situate French and Spanish in a typology of focus realization. This question is addressed in the next section.

\section{A typology of focus realization}
As mentioned in Chapter 1, previous accounts of focus and its realization tend to adopt a strict characterization of languages and their way to express IS-functions (e.g., \cite{lambrecht1996information}, see Chapter 1). Under such a typology, French is said to use constructional strategies, while Spanish employs word order alternations. \\
These conclusions were often based on theoretical studies of the two languages, which, as mentioned in Chapter 2, tend to identify one single means to express focus in each language, thus excluding variability (e.g., Minimalist or OT-approaches and the VS order in Spanish). Later, experimental studies have contributed to weaken these claims, and have shown that speakers have at their disposal several means. The corpus study conducted in this book has contributed to support this view, adding elliptical structures and their frequency to the possible strategies that had already been identified in the literature. \\
As indicated by the results of the corpus query in Chapter 3, certain strategies for expressing focus are more frequent than others. However, this does not imply that speakers lack alternative means to convey the same function. Classifying Spanish as a ``word order language” in contrast to an ``intonation language” would fail to accurately represent the language's complexity. While the corpus study revealed that postverbal subjects are the predominant means of expressing subject focus in Spanish, other strategies — such as prosodic and constructional approaches — are also employed. \\
Similarly, characterizing French solely as a language that utilizes constructional strategies for focus expression overlooks the phenomenon's complexity. Although the corpus results indicate that clefts are the most frequent means of marking focus in French, speakers also use prosodic strategies to convey focus.
Thus, even when defining a typology based on frequencies rather than categorical judgments derived from personal introspection, it remains essential to conduct extensive investigations across various corpora of spontaneous speech. This would yield a more realistic portrayal of the language and provide a stronger empirical foundation for analysis.\\
The extensive empirical research conducted over the past decade has shown that certain focus-marking strategies, although attested, occur less frequently than others. Conversely, some strategies are not only more common but also fulfil multiple information-structural functions, thereby challenging the notion of a strict one-to-one correspondence between form and function. A natural follow-up question, then, is whether this previously assumed one-to-one mapping resulted merely from a lack of empirical investigation, or whether certain strategies and word order patterns have in fact become more or less permissive over time. This issue is addressed in the following section.

\section{Syntactic variation and language change}
Variationist sociolinguistics has established that there is a robust connection between the variation found in any community of speakers at a given point in time and the long-term processes of change studied by historical linguists. He showed that synchronic variation is very often the root of diachronic change. By conducting longitudinal studies on different generations of speakers in New York, he demonstrated that synchronic variation is a reflex of long-term trends \parencite{weinreich1968empirical}.\\
As is often the case in variationist sociolinguistics, numerous studies have explored the relationship between synchronic variation and diachronic change in phonological variables. Differently from the phonetic/phonological variables, however, syntax is generally more impermeable to change, and syntactic shifts tend to occur in a much larger span of time.
A type of language change that has received considerably less attention is the one investigating shifts at the level of the information structure of a given a language, or, more precisely, changes occurring at the interface between syntax and information structure. Intuitively, variation and change at the phonetic/phonological level is easier to grasp than changes in the pragmatic constraints of a language, which are, by definition, more subtle.\\
Nevertheless, there have been some studies investigating the changes in the mapping of IS onto syntax, especially in the passage from Latin to Modern Romance languages. For instance, Medieval Romance varieties are generally claimed to display a V2 word order, featuring the fronting of one or more constituents to the left of the inflected verb (\cite{beninca2006detailed, ledgeway2012latin, salvi2011some, wolfe2015medieval, wolfe2024conspiracy}, \textit{inter alia}). This position was generally reserved for different IS-functions.
In Modern Romance languages, it is possible to observe that certain word orders are not available in equal measure for all the members of the family. \\
\textcite{lahousse2012word}, for instance, investigate the availability of the VOS order in Spanish, Italian and French. 
To account for the fact that this pattern is almost absent in Modern French, the authors postulate an  hypothesis in terms of progressive grammaticalization with respect to the interface between grammatical structure and information structure. They provide evidence for a continuum from Spanish → Italian → French, where Spanish is the least grammaticalized, and French the most grammaticalized language. Additionally, data provided by \textcite{dufter2008explaining} shows that there is a significant increase in the frequency of clefts from Latin to Modern French. The virtual loss of VOS in French, thus, seems to be compensated by the increasing use of these constructions in Modern French. \\
These studies illustrate how changes at the interface between information structure and syntax have influenced the distinctions among Modern Romance languages today. They further support the argument that language change can only be fully understood by considering pragmatic constraints and including these structures within the same envelope of variation.
Additionally, as noted by \textcite{lehmann2008grammaticalization}, ``on the one hand, information structure is present in the source constructions that undergo grammaticalization, and may direct their course. On the other hand, information structure is itself susceptible to grammaticalization” \parencite[1]{lehmann2008grammaticalization}. It is possible that certain word orders expressing multiple IS-functions become restricted to a single pragmatic function.\\
For example, the XPV order was available in many Medieval Romance varieties and is said to express various information-structural functions. Modern Italian has been largely claimed to have limited this position — more generally, the fronting of a constituent — to contrastive or corrective foci (However, see \cite{cruschina2021} on how gradualness in the notion of contrast can account for different subject positions in Italian). Conversely, certain constructions or word orders that were initially restricted to specific IS-functions may become more permissive in terms of pragmatic interpretation, allowing for use in discourse contexts different from their original ones. \textcite{caloi2018answering}, for instance, conducted experiments on Heritage Italian showing that even non-contrastive focused subjects, when accompanying transitive or unergative verbs, may appear in SV rather than VS order. These findings suggest two main points: first, they indicate that a change may be ongoing in this variety of Italian; second, they once again highlight the crucial role of argument structure in the realization of focus. \\
\textit{Il y a} clefts have traditionally been described as presentational or event-reporting constructions, typically associated with an all-focus interpretation \parencite{lambrecht1996information}. However, recent corpus studies have shown that these constructions can also express narrowly-focused subjects \parencite{karssenberg2017oiseaux, Karssenberg2018}. The results obtained from the \textit{sgs}corpus further confirms that \textit{il y a} clefts exhibiting a focus-background partition are not uncommon, supporting the claim that these structures can serve more than one IS-function.
Two factors may account for this phenomenon: (i) earlier studies on \textit{il y a} clefts were primarily based on introspection rather than actual spoken data, whereas more recent studies rely on corpora, offering a fuller perspective on the functions conveyed by this cleft format; (ii) there may have been a change in the IS-functions these constructions express. It is plausible that \textit{il y a} clefts were initially used predominantly in all-focus contexts, and that over time their use has expanded to include narrow-focus cases.\\
Detecting such change in progress is challenging. A longitudinal study — whether in real or apparent time — would provide clearer insights into whether these constructions have undergone shifts in their information-structural roles. This would, however, require speech data collected across multiple generations of speakers or at different time points.

\section{Limitations of the study and next steps}

While the studies presented in this work provide valuable insights into the interaction between information structure, syntax, and argument structure, several limitations must be acknowledged. These limitations concern both methodological and empirical aspects of the research, such as the scope of the data, the representativeness of the corpus and experimental samples, and the interpretative constraints inherent to the chosen analytical framework. Recognizing these boundaries is essential, as it helps to situate the findings within their proper empirical and theoretical context and to identify directions for future research. 

\subsection{Caveats of the study}

The annotation of focus in this study relied on the Question under Discussion (QUD) approach, which provides a principled and replicable way to operationalize how a sentence is partitioned into background and focus. While the QUD framework offers clear advantages for systematically identifying the information-structural function of the utterances, it also presents certain challenges. For instance, determining the precise QUD can be difficult in cases where multiple, equally plausible questions could underlie the same utterance, or where the discourse context does not explicitly indicate the speaker's intended question. Additionally, the method requires interpretation of the pragmatic context, which introduces some subjectivity in the annotation process. These factors suggest that, although the QUD approach is robust, complementary methods may further strengthen the reliability of focus annotation in future studies.\\
Furthermore, in the present annotation scheme, three levels of QUD were distinguished: immediate QUD, distant QUD, and implicit QUD. While this categorization allowed for a systematic analysis of focus, the implicit QUD category conflates two potentially distinct dimensions: immediate-implicit QUD, where the underlying question is closely related to the preceding discourse, and distant-implicit QUD, where the inferred question is more removed from the local context. By merging these two subtypes into a single category, the analysis may have obscured subtle differences in how focus interacts with discourse structure. Future research could benefit from a more fine-grained annotation that separates these dimensions, potentially yielding more detailed insights into the interaction between focus and discourse context.\\
A further limitation relates to the annotation of information status. In the present study, the information status of the subject referent was annotated on the basis of textual givenness, without distinguishing pragmatic givenness or introducing a separate category for cases of bridging. This methodological choice ensured consistency and replicability in annotation, but it may have oversimplified the pragmatic nuances of discourse referents. A more fine-grained distinction—taking into account inferable or bridging referents—could potentially yield different patterns, especially in contexts where the accessibility of information is determined by shared world knowledge rather than explicit textual mention. Future research could therefore benefit from adopting a more detailed annotation scheme that integrates pragmatic dimensions of givenness alongside textual ones.\\
A further limitation concerns the nature of the data used for the analysis. The corpus was based on semi-spontaneous speech elicited through a specific task in which participants were asked to solve a fictional murder case by answering questions. While this method allowed for the controlled elicitation of target structures and ensured comparability across speakers, it does not fully guarantee the spontaneity of the data. Participants were aware of the experimental context and of the need to provide informative answers, which may have influenced their linguistic choices, particularly in relation to focus marking and information packaging. Consequently, some of the patterns observed might reflect task-related effects rather than naturally occurring discourse behaviour. Future research should therefore complement these findings with data from more spontaneous conversational contexts.\\
Finally, beyond the issue of the spontaneity of the data lies the problem of the infrequency of certain constructions. While both the Spanish and French datasets demonstrate the availability of several variants to express the same function, some of these variants are highly infrequent. In the Spanish sub-corpus, for example, speakers produced only 33 \textit{in situ} foci, 29 clefts, and 21 reduced clefts out of 352 items. Similarly, in the French corpus, \textit{in situs}, fragments, and dislocations appear relatively rare, with just 41, 31, and 7 occurrences, respectively.
This raises the question of whether these infrequent variants are genuinely comparable to the more commonly used structures, such as postverbal subjects, French clefts, or Spanish fragments. This question cannot be easily answered by examining frequency data alone. The phenomenon of subject focus is far more intricate and multifaceted. To provide a comprehensive understanding of the contexts in which it occurs and the resulting syntactic structures, it is essential to go beyond purely linguistic and discursive factors, incorporating extra-linguistic variables as well.\\
For instance, the study conducted in this book does not consider the role of register. The formal/informal dichotomy is known to significantly influence various linguistic phenomena, and the mapping of information structure onto syntax is no exception. An investigation of subject focus across different registers would likely account for the scarcity of certain structures. \\
In French, for example, the divergence between formal and informal language is said to be more pronounced than in other languages (see \cite{blanche2006linguistic, katz2007teaching, hamlaoui2011role}, \textit{inter alia}). Standard French (or formal French) refers to the highly codified variety learned later in life, typically in formal contexts and under the influence of literacy. While dominant in written language, Standard French also encompasses the spoken language of educated speakers in formal settings, such as media broadcasts or public speeches.
Conversely, Colloquial French (informal French) refers to the variety used by most native speakers in everyday interactions, acquired early in life in naturalistic settings. Though primarily associated with spoken language, Colloquial French also appears in certain informal written contexts. \\
The differences between these varieties have even raised the question of whether speakers possess two distinct grammars. Variationist approaches, following the work of \textcite{labov1972sociolinguistic}, argue against this, suggesting that natural languages are inherently variable. Other scholars, however, propose a diglossic explanation, suggesting that speakers may have two overlapping but distinct grammars that coexist \parencite{rowlett2011syntactic}.
In her investigation of subject and object focus, \textcite{destruel2016focus} compares the use of formal and informal registers in French, revealing that while clefts are the most common strategy for expressing subject focus in the colloquial variety, speakers in formal settings strongly preferred the \textit{in situ} strategy. This demonstrates a significant effect of register, even when examining information-structural functions.\\
Additionally, social factors have not been considered in this study. Variation between different linguistic forms often mirrors social differences among speakers, such as education level or social status. Incorporating social variables could, therefore, provide further insight into why certain structures are more frequent than others.

\subsection{For further research}
A puzzling result that emerged from the data is the marked difference in the productivity of reduced clefts and fragments between French and Spanish. In French, reduced clefts are highly productive (n=137), while they are infrequent in Spanish (n=21). Conversely, fragments are common in Spanish (n=102) but scarce in French (n=31). To explain this discrepancy, I have proposed the ``copy” hypothesis (see Chapter 3, section 3.7). According to this hypothesis, when speakers aim to produce an elliptical structure for discourse-related reasons, they tend to ``copy” the focused part of the full structure that is more prevalent in their language, simply eliding the background (BG) proposition. Since clefts are the most frequent means of expressing subject focus in spoken French, French speakers producing elliptical structures would likely ``copy” the clefted element and omit the BG proposition contained in the relative clause. \\
Similarly, in Spanish, where postverbal subjects are the most common structure, speakers would replicate the focused element and delete the remaining part of the utterance containing discourse-old information. As a result, the elliptical structure in Spanish manifests as a fragment.
Clearly, this hypothesis requires empirical testing. Ideally, the same study should be replicated using a different dataset or by focusing on other constituents, such as object focus. Additionally, since elliptical constructions are common across many, if not all, languages, this hypothesis could be tested cross-linguistically. A comparative study on focus and the alternation between full and elliptical structures across languages could yield valuable insights into the patterns and properties of elliptical answers, revealing both differences and commonalities among languages.\\
Moreover, the ``copy” hypothesis does not fully explain the occurrence of fragments in French or reduced clefts in Spanish. In other words, it fails to account for the presence of alternative elliptical strategies within the same language. Although these structures appear with much lower frequency, Spanish speakers still produced some reduced clefts, and French speakers some fragments. As mentioned earlier, social or stylistic factors—which could provide further insight—were not considered in this study. At least for French, another possible explanation might be linked to the type of verb in the preceding QUD (which is elided in the answer). Some linguistic contexts may render a reduced cleft infelicitous, while a fragment would be more appropriate, and conversely, other contexts might favour the use of a reduced cleft over a fragment. Consider the following examples: 

  
\ea \label{ex92}
\ea \textit{Qui veut du gâteau?}\\
 ‘Who wants cake?'
\exi{a.} \ [\textit{Moi}]\textsubscript{\textsc{foc}}\normalsize!\\
‘Me!'
\exi{b\small{$'$}.} \#\textit{C'est} [\textit{moi}]\textsubscript{\textsc{foc}}\normalsize!\\
‘It is me!'
\z \z 

\ea \label{ex93}
\ea \textit{C'est qui dans la photo?}\\
 ‘Who is this in the picture?'
\exi{b.} \#[\textit{Moi}]\\textsc{foc}\normalsize!\\
‘Me!'
\exi{b\small{$'$}.}  \textit{C'est} [\textit{moi}]\ \textsc{foc}\normalsize!\\
‘It is me!'
\z \z

In example (\ref{ex92}), a fragment seems to be a more felicitous answer than a reduced cleft. Conversely, in example (\ref{ex93}), a reduced cleft appears to be a better option. While it is not yet clear what determines these preferences, it seems that certain contexts impose more constraints than others, and these constraints may be related to the type of verb (and possibly the thematic roles) in the question.
It is important to emphasize that answering with a reduced cleft in the context of (\ref{ex92}) or using a fragment in (\ref{ex93}) would not result in ungrammaticality. Rather, it reflects probabilistic choices and speaker preferences. What exactly drives these preferences remains an open question and offers a promising direction for future research.




